%==============================================================================
% BAB III: LANDASAN TEORI
% [Constraints]:
% - Sumber wajib ilmiah: buku, jurnal, prosiding, skripsi/TA, web resmi terpercaya
% - DILARANG menggunakan Wikipedia, blogspot, koran, atau portal berita non-ilmiah
%==============================================================================
\chapter{Landasan Teori}

\pet{Pada bagian ini dijelaskan mengenai konsep dan dasar teori yang digunakan dalam implementasi sistem, aplikasi, atau prototype yang dikembangkan.}

\section{Dasar Teori}
\pet{Berisi teori/konsep yang berkaitan/digunakan dalam kerja praktek. Gunakanlah data melalui buku/jurnal referensi, publikasi tugas akhir, penelitian, buku, dan informasi web yang dapat dipertanggungjawabkan, \petb{hindari penggunaan dasar teori melalui tautan Wikipedia, surat kabar, atau portal berita, yang dapat memiliki isi yang tidak bersifat fakta.}}

\section{Teori I}
\pet{Deskripsikan mengenai teori/konsep yang berkaitan/digunakan/menjadi acuan dalam penelitian. Kemudian berikan pembahasan sederhana mengenai penggunaannya di dalam laporan PKL yang Anda kerjakan.}

\section{Teori II}
\pet{Deskripsikan mengenai teori / konsep yang berkaitan / digunakan / menjadi acuan dalam penelitian. Kemudian berikan pembahasan sederhana mengenai penggunaannya di dalam tugas akhir yang Anda kerjakan.}

\subsection{Sub Teori I}
\pet{Sub Teori dapat dituliskan dengan level penomoran 3 \& style “Heading 3”}
